\section{Custom CNN Architectures}
\subsection{Model A}
Four standard $3\times3$ convolutions use channel widths 24, 48, 96 and 128.
Each is followed by batch normalization, ReLU and $2\times2$ max-pooling.
Padding preserves convolutional spatial size. Global average pooling and a
small linear head avoid a large flattened dense layer. Batch normalization
provides learned channel scaling and stabilizes intermediate feature scales.
Table~\ref{tab:model_a_layers} gives every layer's input/output shape, kernel,
stride, padding, activation and parameter count. The total is \AParams{};
all parameters are trainable. The linear head outputs ten unnormalized logits
for cross-entropy loss.

\subsection{Model B}
The stem, channel widths, pooling schedule and head match Model A. Each later
standard convolution is replaced by a depthwise $3\times3$ convolution followed
by a pointwise $1\times1$ convolution. Depthwise groups equal input channels.
Each component has batch normalization and ReLU. These extra operations mean
this is not a strict single-operator ablation. Table~\ref{tab:model_b_layers}
identifies depthwise and pointwise operations separately. Channel counts are
the second dimension of each NCHW shape, and activations are shown as separate
rows. Model B is built from scratch using a reusable
\texttt{DepthwiseSeparableConv} block.

The executed count is \BParams{} trainable parameters, leaving
\BMargin{} below the 100,000 limit. The constructor and verification script
check this programmatically:
\begin{verbatim}
count = sum(p.numel() for p in model.parameters()
            if p.requires_grad)
assert count < 100000
\end{verbatim}
The sum of the layer counts in Table~\ref{tab:model_b_layers} independently
agrees with this result.

\subsection{Depthwise Separable Convolution}
Depthwise filtering processes channels independently; pointwise convolution
learns cross-channel combinations \cite{howard2017}. For output spatial size
$H_o\times W_o$, multiplying the weight expressions below by $H_oW_o$ gives
the corresponding convolution MAC counts for one image.
Here a MAC is one multiply--accumulate operation. The factorization reduces
the number of weights because a separate spatial kernel is no longer learned
for every input--output channel pair. It also restricts the possible filters,
so equal accuracy is not guaranteed.
\subsection{Parameter Calculation}
With convolution biases disabled,
\begin{align}
P_{\mathrm{standard}} &= K_hK_wC_{\mathrm{in}}C_{\mathrm{out}},\label{eq:standard}\\
P_{\mathrm{DSC}} &= K_hK_wC_{\mathrm{in}} + C_{\mathrm{in}}C_{\mathrm{out}}.\label{eq:dsc}
\end{align}
In Equations~\eqref{eq:standard} and~\eqref{eq:dsc}, $K_h,K_w$ are kernel
dimensions and $C_{\mathrm{in}},C_{\mathrm{out}}$ are channel counts. Dividing
the two expressions gives
\begin{equation}
\frac{P_{\mathrm{DSC}}}{P_{\mathrm{standard}}}
=\frac{1}{C_{\mathrm{out}}}+\frac{1}{K_hK_w},
\qquad
\text{reduction}=100\left(1-\frac{P_{\mathrm{DSC}}}{P_{\mathrm{standard}}}\right)\%.
\label{eq:reduction}
\end{equation}
Table~\ref{tab:savings} applies Equation~\eqref{eq:reduction} to the three
replaced convolutions. Batch normalization adds two
trainable scalars per output channel; running statistics are non-trainable
buffers. The classification head includes bias. Thus whole-network counts
include more than the convolution weights in this table.
% Generated from saved experiment evidence; do not edit values.
\begin{table}[H]\centering\small
\caption{Convolution weight counts excluding bias and normalization; reduction compares depthwise plus pointwise with standard convolution.}\label{tab:savings}
\begin{tabular}{lrrrr}
\toprule
Channels & Standard & Depthwise & Pointwise & Reduction (\%) \\ \midrule
24 to 48 & 10368 & 216 & 1152 & 86.81 \\
48 to 96 & 41472 & 432 & 4608 & 87.85 \\
96 to 128 & 110592 & 864 & 12288 & 88.11 \\
\bottomrule
\end{tabular}
\end{table}

Summing the executed architecture tables gives \AConvParams{} convolution
weights, \ABnParams{} normalization parameters and \AHeadParams{} head
parameters for Model A, totalling \AParams{}. The corresponding Model B
counts are \BConvParams{}, \BBnParams{} and \BHeadParams{}, totalling
\BParams{}. Whole-network savings differ from an individual convolution's
savings because the stem and head are retained and Model B adds normalization.

\subsection{Hardware-Aware Activation Selection}
ReLU computes $\max(0,x)$ using a simple comparison and avoids exponentials.
It is widely supported by inference kernels and is compatible with common
integer inference workflows. No quantization or activation-specific energy
measurement is performed here, so hardware savings are not claimed empirically.
ReLU is the deliberate choice for both custom networks. The pretrained
architectures retain their original activations: MobileNetV2 uses ReLU6,
whereas EfficientNet-B0 uses SiLU and sigmoid operations. They are not modified
to match the custom networks.
